\subsubsection{QOJ837 Giant Penguin}

\problemmeta{03/13/2026}{}{Centroid Decomposition, DFS Tree}

\paragraph{Statement}
Given an undirected graph with $n$ nodes and $m$ edges, $n \le 10^5, m \le 2\times10^5$. Each node is on at most $k \le 10$ cycles. You will perform $q \le 2\times10^5$ operations.
\begin{enumerate}
    \item Mark a node $u$.
    \item Find the shortest distance from any marked node to node $u$.
\end{enumerate}

\paragraph{Editorial} Consider this problem, when the graph is a tree, how would you do it.
The obvious solution would be doing a centroid decomposition, like the \textit{Xenia and Tree} problem.
But because of this problem is a grpah, we can't do it.

Consider a centroid node $c$, there will be at most $k$ edges connecting different subtrees (since $c$ can have at most $k$ cycles).
We can just do a BFS from one of the end of that edge, and find the shortest path from it to the others.
We call that node \textit{special}.

When we do query Type 1, we can traverse up its centroid parents, and also each parent's \textit{special} node.
We will update by $\text{ans}[c][sp] = \text{min}(\text{ans}[c][sp], \text{dist}[c][sp][u])$.
For Type 2, we can do similarly, $\text{ANS} = \text{min}(\text{ANS}, \text{ans}[c][sp] + \text{dist}[c][sp][u])$.

But since $\text{dist}[c][sp][u]$ size is $\mathcal{O}(n^2k)$, we have to change $c$ to $\text{dep}[c]$. This works because there is only one $c$ that has depth = $\text{dep}[c]$ and is the parent of $u$.

\codelink{CNHW/2026/2.cpp}